\subsection{单元素}

定义 2。若存在 \(h\)，使  为 \(\mathfrak{g}\) 中的一个 \(x, y \in \mathfrak{g}\)\((x, h, y)\)\(\mathfrak{sl}_2\) 三元组，则称 \(\mathfrak{g}\) 的元素 h 为单元素。

我们也称 \(h\) 为 \(\mathfrak{sl}_2\) 三元组 \((x, h, y)\) 的单元素。

命题 4。令 \(h\) 为 \(\mathfrak{g}\) 的非零元素。则 \(h\) 是单元素，当且仅当存在 \(x \in \mathfrak{g}\)，使 \([h, x] = 2x\) 且 \(h \in (\operatorname{ad} x)(\mathfrak{g})\)。

必要性显然。充分性由引理6 得到。

命题 5。假设 g 是可分裂半单李代数。令 h 为 g 的分裂嘉当子代数，R 为 \((\mathfrak{g},\mathfrak{h})\) 的根的集合，B 为 R 的一个基。令 h 为属于 h 的 g 的一个单元素。则 h 在 \(\operatorname{Aut}_{e}(\mathfrak{g},\mathfrak{h})\) 下与 h 中满足对所有 \(h'\) 都有 \(\alpha(h') \in \{0,1,2\}\) 的元素 \(\alpha \in B\) 共轭。

\(\mathrm{ad}_{\mathfrak{g}}h\) 的特征值属于 \(\mathbf{Z}\)（§1，第2号，命题2的推论）。因此 \(h \in \mathfrak{h}_{\mathbf{Q}}\)。存在 \(w\) 的 Weyl 群中的元素 \((\mathfrak{g}, \mathfrak{h})\)，使对所有 \(\alpha(wh) \geq 0\) 都有 \(\alpha \in \mathrm{B}\)（第 VI 章第1节第5号定理2 (i)）。根据 §2 第2号定理2的推论，我们只需考虑对所有 \(\alpha(h) \in \mathbf{N}\) 都有 \(\alpha \in \mathrm{B}\) 的情形。令 \(\mathrm{R}_{+}\) 为关于 B 的正根集合，\(\mathrm{R}_{-} = -\mathrm{R}_{+}\)。\(\mathfrak{sl}_2\) 中存在形如 \(\mathfrak{g}\) 的 \((x, h, y)\) 三元组。令 T 为满足 \(\beta\) 的根 \(\beta(h) = -2\) 所成的集合。则 \(\mathrm{T} \subset \mathrm{R}_{-}\)，且 \(y \in \sum_{\beta \in \mathrm{T}} \mathfrak{g}^{\beta}\)。假设存在 \(\alpha \in \mathrm{B}\) 使 \(\alpha(h) > 2\)。对所有 \(\beta \in \mathrm{T}\)，有 \((\alpha + \beta)(h) > 0\)，所以 \(\alpha + \beta \notin \mathrm{R}_{-}\) 且 \(\alpha + \beta \neq 0\)；另一方面，由于 \(\beta \in \mathrm{R}_{-}\) 且 \(\alpha \in \mathrm{B}\)，有 \(\alpha + \beta \notin \mathrm{R}_{+}\)；因而 \(\alpha + \beta \notin \mathrm{R} \cup \{0\}\)，所以 \([\mathfrak{g}^{\alpha}, \mathfrak{g}^{\beta}] = 0\)。于是 \([y, \mathfrak{g}^{\alpha}] = 0\)。但 \(\mathrm{ad}_{\mathfrak{g}}y|\mathfrak{g}^{\alpha}\) 是单射，因为 \(\alpha(h) > 0\)（§1，第2号，命题2的推论）。这个矛盾证明了对所有 \(\alpha(h) \leq 2\) 都有 \(\alpha \in \mathrm{B}\)。

推论。若 k 是代数闭域且 \(k\)\(\mathfrak{g}\) 是秩为 \(l\) 的半单李代数，则相对于 ，\(\mathfrak{g}\) 的单元素共轭类至多有 \(\operatorname{Aut}_e(\mathfrak{g})\)\(3^l\) 个。

事实上，\(\mathfrak{g}\) 的每个半单元素都在 \(\mathrm{Aut}_e(\mathfrak{g})\) 下与 \(\mathfrak{h}\) 中的一个元素共轭。

引理 8。假设 k 是代数闭域且 \(k\)\(\mathfrak{g}\) 是半单的。令 \(h\) 为 \(\mathfrak{g}\) 的半单元素，且 \(\operatorname{ad} h\) 的特征值都是有理数。令 \(\mathfrak{g}^0 = \operatorname{Ker}(\operatorname{ad} h)\)，\(\mathfrak{g}^2 = \operatorname{Ker}(\operatorname{ad} h - 2)\)。令 \(G_h\) 为保持 h 不变的 \(\mathfrak{g}\) 的初等自同构所成的集合。令 \(h\)\(x \in \mathfrak{g}^2\) 满足 \([x, \mathfrak{g}^0] = \mathfrak{g}^2\)。则 \(G_h x\) 包含 \(\mathfrak{g}^2\) 的一个在 Zariski 拓扑下稠密且开的子集。

令 \(\mathfrak{h}\) 为 \(\mathfrak{g}^0\) 的一个嘉当子代数。这是包含 h 的 \(\mathfrak{g}\) 的一个嘉当子代数（第 VII 章第2节第3号命题10）。我们有 \(h\)\(h \in \mathfrak{h}_{\mathbf{Q}}\)。令 R 为 \((\mathfrak{g}, \mathfrak{h})\) 的根系，Q 为根权格。存在 R 的一个基 B，使对所有  都有 \(\alpha(h) \geq 0\)。\(\alpha \in B\)

令 \(\mathrm{U}\) 为满足对所有  都有  的 \(z \in \mathfrak{h}\) 所成的集合。令 \(\alpha(z) \neq 0\)\(\alpha \in \mathrm{B}\)\((H_{\alpha}')_{\alpha \in \mathrm{B}}\) 为与 B 对偶的 \(\mathfrak{h}\) 的基。若 \(\mathrm{B}\)\(z \in \mathrm{U}\)，则存在从 \(\mathrm{Q}\) 到 \(k^*\) 的同态，它将任意 \(\gamma \in \mathrm{Q}\) 映为 \(\prod_{\alpha \in \mathrm{B}} \alpha(z)^{\gamma(H_{\alpha}')}\)。根据 §5 命题2和4，在向量空间  上的自同态 \(\varphi(z)\) 在 \(\mathfrak{g}\)\(\mathfrak{g}^{\gamma}\) 上诱导比值为 \(\prod_{\alpha \in \mathrm{B}} \alpha(z)^{\gamma(H_{\alpha}')}\) 的伸缩，因此是 \(\mathfrak{g}\) 的初等自同构，并且显然属于 \(G_h\)。

令 \(s \in \mathfrak{h}\)。若 \(\gamma \in \mathrm{R}\) 满足 \(\mathfrak{g}^{\gamma} \cap \mathfrak{g}^{2} \neq 0\)，

\[
2 = \gamma (h) = \gamma \left(\sum_ {\alpha \in \mathrm{B}} \alpha (h) H _ {\alpha} ^ {\prime}\right) = \sum_ {\alpha \in \mathrm{B}} \alpha (h) \gamma (H _ {\alpha} ^ {\prime});
\]

由于对所有 \(\alpha(h) \geq 0\) 都有 \(\alpha \in \mathrm{B}\)，且 \(\gamma(H_{\alpha}^{\prime})\) 是全都 \(\geq 0\) 或全都 \(\leq 0\) 的整数，故对所有 \(\gamma(H_{\alpha}^{\prime}) \in \mathbf{N}\) 都有 \(\alpha \in \mathrm{B}\)。因此，我们可以（对 \(z \in \mathfrak{h}\)）考虑向量空间 \(\psi(z)\) 的自同态 \(\mathfrak{g}^2\)，它在 \(\mathfrak{g}^{\gamma} \cap \mathfrak{g}^2\) 上诱导比值为 \(\prod_{\alpha \in \mathrm{B}} \alpha(z)^{\gamma(H_{\alpha}^{\prime})}\) 的伸缩。映射 \(z \mapsto \psi(z)\) 从 \(\mathfrak{h}\) 到 \(\operatorname{End}(\mathfrak{g}^2)\) 是多项式的。对 \(z \in \mathrm{U}\)，有 \(\psi(z) = \varphi(z)|\mathfrak{g}^2\)。

令 \(\gamma_1, \ldots, \gamma_r\)\((\mathfrak{g}, \mathfrak{h})\)\(h\)\(y_1 \in \mathfrak{g}^{\gamma_1}, \ldots, y_r \in \mathfrak{g}^{\gamma_r}\)\(e^{\mathrm{ad} y_1} \ldots e^{\mathrm{ad} y_r} \in \mathrm{G}_h\)\(\rho\) 为在 h 上消失的  的互不相同的根。若 ，则 。于是可定义从 \(\mathfrak{h} \times \mathfrak{g}^{\gamma_1} \times \cdots \times \mathfrak{g}^{\gamma_r}\) 到 \(\mathfrak{g}^2\) 的映射 ，置

\[
\rho (z, y _ {1}, \dots , y _ {r}) = \psi (z) e ^ {\operatorname{ad} y _ {1}} \dots e ^ {\operatorname{ad} y _ {r}} x
\]

对 \(z \in \mathfrak{h}\)、\(y_1 \in \mathfrak{g}^{\gamma_1}, \ldots, y_r \in \mathfrak{g}^{\gamma_r}\)，此映射是多项式的，且 \(\rho(\mathrm{U}, \mathfrak{g}^{\gamma_1}, \ldots, \mathfrak{g}^{\gamma_r}) \subset \mathrm{G}_h x\)。根据第 VII 章附录 I 命题3和4，只需证明 \(\rho\) 的切线性映射在某一点满射。

现在令 T 为 \(z \mapsto \psi(z)\) 在 \(h_{0} = \sum_{\alpha \in B} H_{\alpha}'\) 处的切线性映射。于是 \(\mathrm{T}(z)\) 是 \(g^{2}\) 的自同态，它在 \(g^{\gamma} \cap g^{2}\) 上诱导比值为

\[
\sum_ {\alpha \in \mathrm{B}} \gamma (H _ {\alpha} ^ {\prime}) \alpha (h _ {0}) ^ {\gamma (H _ {\alpha} ^ {\prime}) - 1} \alpha (z) \prod_ {\beta \in \mathrm{B}, \beta \neq \alpha} \beta (h _ {0}) ^ {\gamma (H _ {\alpha} ^ {\prime})} = \sum_ {\alpha \in \mathrm{B}} \gamma (H _ {\alpha} ^ {\prime}) \alpha (z) = \gamma (z).
\]

因此，\(z \mapsto \rho(z,0,\ldots,0)\) 在 \(h_{0}\) 处的切线性映射是 \(z \mapsto [z,x]\)，其像为 \([x,h]\)。映射 \(y_{1} \mapsto \rho(h_{0},y_{1},0,\ldots,0)\) 在 0 处的切线性映射是 \(y_{1} \mapsto \psi(h_{0})[y_{1},x]\)；此映射的像为 \(\psi(h_{0})[x,\mathfrak{g}^{\gamma_{1}}] = [x,\mathfrak{g}^{\gamma_{1}}]\)。类似地，映射 \(y_{i} \mapsto \rho(h_{0},0,\ldots,0,y_{i},0,\ldots,0)\) 在 0 处的切线性映射的像为 \([x,\mathfrak{g}^{\gamma_{i}}]\)。最后，\(\rho\) 在 \((h_{0},0,\ldots,0)\) 处的切线性映射的像为

\[
[ x, \mathfrak {h} + \mathfrak {g} ^ {\gamma_ {1}} + \dots + \mathfrak {g} ^ {\gamma_ {r}} ] = [ x, \mathfrak {g} ^ {0} ] = \mathfrak {g} ^ {2}.\tag{Q.E.D.}
\]

*群 \(\mathrm{G}_h\) 是一个代数群，其李代数为 ad \(\mathfrak{g}^0_*\)

命题 6。假设 k 是代数闭域且 \(k\)\(\mathfrak{g}\)\(G\)\(\mathfrak{g}\) 是半单的。令 G 为包含 \(\operatorname{Aut}_e(\mathfrak{g})\) 的 \((x, h, y)\)\((x', h', y')\) 的自同构群。令  和  为  中的 \(\mathfrak{sl}_2\) 三元组。以下条件等价：\(\mathfrak{g}\)

\begin{enumerate}
\def\labelenumi{(\roman{enumi})}
\item
  h 与 \(h'\) G-共轭；
\item
  \((x,h,y)\) 与 \((x^{\prime},h^{\prime},y^{\prime})\) G-共轭。
\end{enumerate}

只需证明蕴含 (i) \(\Longrightarrow\) (ii)，且可立即归约到 \(h = h'\) 的情形。引入引理8中的 \(g^{2}\) 和 \(G_{h}\)。由 §1，第2号，命题2的推论，有 \(x \in g^{2}\) 且 \([x, g^{0}] = g^{2}\)。因此，\(G_{h}x\) 包含 \(g^{2}\) 的一个在 Zariski 拓扑下稠密且开的子集，\(G_{h}x'\) 也如此。所以存在 \(a \in G_{h}\)，使 \(a(x) = x'\)。我们有 \(a(h) = h\)，因而 \(a(y) = y'\)（第1号，引理1）。

推论 1。将任意 \(\mathfrak{sl}_2\) 三元组与其单元素关联的映射在取商后，定义了从 \(\mathfrak{sl}_2\) 三元组的 G-共轭类到单元素的 G-共轭类的一个双射。

推论 2。若 \(\operatorname{rk}(\mathfrak{g}) = l\)，则相对于 \(\operatorname{Aut}_e(\mathfrak{g})\)，\(\mathfrak{g}\) 的非零幂零元素的共轭类数至多为 \(3^l\)。

这由推论1、命题2的推论以及命题5的推论得到。

推论 3。若 \(\mathrm{rk}(\mathfrak{g}) = l\)，则相对于 \(\mathrm{Aut}_e(\mathfrak{g})\)，同构于 \(\mathfrak{g}\) 的 \(\mathfrak{sl}(2,k)\) 的子代数的共轭类数至多为 \(3^l\)。

这由推论1、命题1以及命题5的推论得到。

